% ProbeRL研究框架黑白流程图
% 使用TikZ绘制，虚线框标明模块，箭头标明关系

\begin{figure}[htbp]
\centering
\begin{tikzpicture}[
    node distance=0.8cm and 1.2cm,
    module/.style={draw, dashed, rounded corners=4pt, inner sep=8pt, fill=white},
    box/.style={draw, solid, rounded corners=2pt, minimum width=2.8cm, minimum height=0.7cm, align=center, font=\small, fill=white},
    arrow/.style={-{Stealth[length=5pt]}, thick},
    biarrow/.style={{Stealth[length=5pt]}-{Stealth[length=5pt]}, thick},
    label/.style={font=\scriptsize, midway},
]

% ===== 攻击侧 =====
\node[box] (threat) {威胁建模};
\node[box, right=of threat] (attack_policy) {GRPO攻击策略\\$\pi_\phi^{\text{atk}}$};
\node[box, right=of attack_policy] (search) {停滞感知\\搜索扩展};
\node[box, below=0.6cm of search] (reward) {步骤级密集奖励\\$\mathcal{R}_{\text{dense}}$};

% 攻击侧虚线框
\node[module, fit=(attack_policy)(search)(reward), label={[anchor=north west, font=\small\bfseries]north west:研究点一：攻击策略}] (atk_mod) {};

% 奖励信号子框
\node[box, below=0.6cm of reward, minimum width=2.8cm, font=\scriptsize] (r_tool) {$\mathcal{R}_{\text{tool}}$: 工具调用偏差};
\node[box, below=0.15cm of r_tool, minimum width=2.8cm, font=\scriptsize] (r_text) {$\mathcal{R}_{\text{text}}$: 文本异常信号};
\node[box, below=0.15cm of r_text, minimum width=2.8cm, font=\scriptsize] (r_env) {$\mathcal{R}_{\text{env}}$: 环境变化信号};
\node[box, below=0.15cm of r_env, minimum width=2.8cm, font=\scriptsize] (r_cover) {$\mathcal{R}_{\text{cover}}$: 覆盖路径新颖性};
\node[box, below=0.15cm of r_cover, minimum width=2.8cm, font=\scriptsize] (r_risk) {$\mathcal{R}_{\text{risk}}$: 风险评分};

% 奖励子框
\node[module, fit=(r_tool)(r_text)(r_env)(r_cover)(r_risk), label={[anchor=north west, font=\scriptsize\bfseries]north west:奖励信号}] (rew_mod) {};

% ===== 防御侧 =====
\node[box, right=2.5cm of attack_policy] (def_policy) {GRPO防御策略\\$\pi_\psi^{\text{def}}$};
\node[box, right=of def_policy] (detect) {语义注入\\检测模型};
\node[box, below=0.6cm of detect] (rule) {启发式规则\\兜底防御};

% 防御侧虚线框
\node[module, fit=(def_policy)(detect)(rule), label={[anchor=north west, font=\small\bfseries]north west:研究点二：防御策略}] (def_mod) {};

% ===== 环境 =====
\node[box, below=2.8cm of attack_policy] (agentdojo) {AgentDojo\\智能体执行环境};

% ===== 箭头关系 =====
% 威胁建模 → 攻击策略
\draw[arrow] (threat) -- (attack_policy);

% 攻击策略 ↔ 搜索扩展
\draw[biarrow] (attack_policy) -- (search) node[label, above] {候选生成/反馈};

% 搜索扩展 → 奖励计算
\draw[arrow] (search) -- (reward) node[label, right] {执行轨迹};

% 奖励 → 各信号
\draw[arrow] (reward) -- (r_tool);

% 攻击策略 → AgentDojo
\draw[arrow] (attack_policy) -- (agentdojo) node[label, left] {注入提示};

% AgentDojo → 奖励
\draw[arrow] (agentdojo.east) -| ([xshift=0.3cm]reward.south east) -- (reward.south east) node[label, right, pos=0.3] {轨迹/状态};

% 攻击侧 → 防御侧: 对抗训练
\draw[biarrow, dashed] (attack_policy) -- (def_policy) node[label, above] {对抗训练};

% 防御策略 → 检测模型
\draw[arrow] (def_policy) -- (detect) node[label, above] {语义过滤};

% 检测模型 → 启发式规则
\draw[arrow] (detect) -- (rule) node[label, right] {漏报兜底};

% AgentDojo → 防御
\draw[arrow] (agentdojo.east) -- ++(1.5,0) |- (def_policy.south) node[label, right, pos=0.25] {攻击样本};

% 防御 → AgentDojo
\draw[arrow] (def_policy.south) -- ++(0,-0.3) -| ([xshift=-0.3cm]agentdojo.east) node[label, below, pos=0.75] {防御部署};

\end{tikzpicture}
\caption{ProbeRL研究框架}
\label{fig:framework}
\end{figure}
